\chapter{Introduction \& Motivation}
\label{ch:introduction}

%==============================================================================
\section{The Shifting Alzheimer's Diagnostic Paradigm}
\label{sec:paradigm}

Alzheimer's disease is not an event but a continuum. The clinical progression runs from
subjective cognitive decline (SCD), in which a person reports memory difficulty that
standardised testing cannot yet confirm, through mild cognitive impairment (MCI), in which
objective deficits appear without loss of functional independence, to Alzheimer's dementia
proper. Pathology precedes symptoms by an estimated one to two decades: amyloid deposition
and tau spread are well advanced by the time a diagnosis is made on clinical grounds.

For most of the field's history, that long prodromal window was of academic interest only,
because nothing could be done with it. Two developments changed that. First, the arrival
of disease-modifying anti-amyloid therapies --- lecanemab and donanemab --- shifted the
clinical question from \emph{whether} a patient has dementia to \emph{whether a patient
without dementia is on a trajectory toward it}, because these agents are indicated in
early symptomatic disease and their benefit depends on being started early. Second, their
prescription requires biomarker confirmation of amyloid pathology, which makes accurate,
low-cost prodromal risk stratification a practical bottleneck rather than a research
aspiration.

The clinically consequential question addressed in this thesis follows directly:
\emph{given a patient who presents with MCI, will they progress to Alzheimer's dementia?}
Framed as a supervised learning problem, this is converter-versus-stable-MCI
classification, and the population it must work on is precisely the population for which
the answer is not already obvious.

%==============================================================================
\section{Resting-State fMRI as a Non-Invasive Window}
\label{sec:rsfmri-motivation}

The biomarkers that currently establish amyloid status are positron emission tomography
and cerebrospinal fluid assay. Both work. Both are poorly suited to population-scale
prodromal screening: PET requires a cyclotron-produced radiotracer, delivers an ionising
radiation dose, and is expensive; CSF assay requires lumbar puncture, which is invasive
and carries a compliance cost that rises steeply in exactly the older, ambivalent
population that early screening targets.

Resting-state functional MRI offers a different trade. It is non-invasive, involves no
ionising radiation, and is acquirable on standard clinical scanners. Its scientific appeal
is that functional \emph{dysconnectivity} --- disruption of the coordinated activity
between brain regions, particularly within the default mode network --- appears
\emph{before} macrostructural atrophy is measurable. If the disease disturbs how regions
communicate before it visibly destroys them, then a functional measure should in principle
carry earlier prognostic information than a structural one.

The cost of that trade is signal quality. The BOLD signal is an indirect, sluggish,
noise-dominated proxy for neural activity, sensitive to head motion, physiological
confounds, and scanner differences. Much of the methodological care documented in
Chapter~\ref{ch:data} exists to prevent those confounds from being learned instead of the
disease.

%==============================================================================
\section{The Actual Research Setting: Deep Learning Under Severe Sample Constraint}
\label{sec:setting}

It is conventional to list small sample size as a limitation in a closing section. In this
work it is the defining condition of the problem, and it is stated here as a premise rather
than an apology, because nearly every result in Chapters~\ref{ch:sota-repair}
and~\ref{ch:results} is a consequence of it.

\paragraph{Dimensionality against sample size.} A single Schaefer-200 functional
connectivity matrix contains $19{,}900$ unique off-diagonal entries. The DELCODE
longitudinal MCI cohort used here contains $167$ subjects, of whom fewer than $70$ are
converters. There are two orders of magnitude more features per subject-visit than there
are subjects. Any model with substantial capacity can fit this data perfectly and learn
nothing generalisable, and distinguishing the two outcomes requires evaluation discipline
that the field does not uniformly practise (Section~\ref{sec:gap-floors}).

\paragraph{Irregular follow-up and attrition.} Subjects are not scanned on a common
schedule, do not contribute equal numbers of visits, and --- critically --- do not drop out
at random. In this cohort converters are followed \emph{longer} than stable subjects
(Section~\ref{sec:gap-attrition}), so sequence length is correlated with the label. A
longitudinal model therefore has access to a shortcut that has nothing to do with brain
function, and demonstrating that it did not take that shortcut requires a specific analysis
rather than an assurance.

\paragraph{Regularity where irregularity was expected.} A finding of this work, reported in
Section~\ref{sec:delta-t-finding}, is that DELCODE's follow-up is in fact highly regular:
90\% of all inter-visit intervals are exactly twelve months, by protocol design. This has
sharp consequences for what a time-aware model can be shown to contribute on this cohort,
and it is one of the reasons external validation on cohorts with genuinely irregular
follow-up is not decoration but a load-bearing part of the argument.

%==============================================================================
\section{Research Objectives \& Core Contributions}
\label{sec:contributions}

This thesis pursues four objectives, and the contributions are stated below at the strength
the evidence supports rather than at the strength the architecture was designed to deliver.

\begin{enumerate}
  \item \textbf{A parameter-lean longitudinal trajectory model.} A recurrent classifier
    (GE-LSTM / GE-GRU) consumes a sequence of per-visit connectome embeddings, with gate
    updates conditioned on the inter-visit interval, and predicts conversion. Its design
    priority is sample efficiency rather than capacity.

  \item \textbf{A pre-registered ablation of self-supervised graph pretraining.} Four arms
    --- no encoder, randomly initialised encoder, pretrained frozen encoder, pretrained
    fine-tuned encoder --- separate the contribution of the reconstruction pretraining, of
    the architecture prior, and of having a graph encoder at all. The interpretation table
    was written before any results existed, so no conclusion could be fitted to the data
    after the fact.

  \item \textbf{A reproducibility audit of a published state-of-the-art baseline.} The
    Brain-TokenGT release is shown to ship with its recurrent core disabled, to diverge when
    that core is enabled, and --- after stabilisation --- to remain \emph{non-reproducible
    at a fixed random seed}, with within-seed variance exceeding between-seed variance
    across 19 runs. This invalidates the error bars of any four-seed summary of that model,
    including one this project itself previously reported.

  \item \textbf{Cross-cohort replication.} ADNI and OASIS-3 were processed through an
    identical preprocessing pipeline so that the central findings can be tested outside the
    cohort that produced them.
\end{enumerate}

\risk{Contribution (1) is the one most exposed by this thesis's own results.}{%
The encoder ablation of Chapter~\ref{ch:results} finds that removing the graph encoder
entirely does not hurt, and the $\Delta t$ ablation of Section~\ref{sec:delta-t-finding}
finds that removing the time conditioning changes the held-out AUC by $0.000$ on DELCODE.
What survives of contribution (1) on this cohort is the \emph{recurrence over visit order},
not the graph encoding and not the time-awareness. Chapter~\ref{ch:discussion} treats this
as the honest headline rather than as a limitation appended to a stronger claim.}

%==============================================================================
\section{Thesis Organization}
\label{sec:organisation}

Chapter~\ref{ch:foundations} establishes the technical background and then --- at greater
length --- sets out the eight questions in this literature that are genuinely unsettled,
because the results that follow are only interpretable as answers to them.
Chapter~\ref{ch:data} documents the DELCODE cohort, the ADNI and OASIS-3 external cohorts,
the preprocessing pipeline, and the leakage-prevention protocol.
Chapter~\ref{ch:methodology} specifies the models. Chapter~\ref{ch:sota-repair} reports the
baseline audit and the reproducibility finding. Chapter~\ref{ch:results} reports the
experimental results and ablations. Chapter~\ref{ch:software} describes the experiment
infrastructure that makes the reproducibility claims checkable.
Chapter~\ref{ch:discussion} interprets, and Chapter~\ref{ch:conclusion} concludes.

Throughout, four kinds of margin annotation mark the structure of the argument: research
gaps in the field, grey areas where the evidence is genuinely contested, points where this
work is defensibly ahead, and points where it is exposed. Appendix~E collects every one of
them into a single register.
